\section{Discussion}
\label{sec:discussion}

\subsection{Principal Findings}

This project set out to determine how much fraud a bank can catch while keeping
the harm to legitimate customers tolerable, and at what cost. The answers are
summarised below.

\paragraph*{Accuracy is the wrong yardstick, and the gap is stark.} A model that
labels every transaction legitimate scores 99.83\% accuracy and catches no
fraud. Our model scores 99.96\% --- a 0.13 percentage-point difference --- and
catches 78.2\% of frauds. Any report that quotes accuracy alone would have
concealed the entire result.

\paragraph*{At 95\% precision, we can catch roughly three quarters of fraud.} The
honest figure is \textbf{73.24\%}, derived from a threshold selected on
out-of-fold training predictions; the test-tuned figure of 78.17\% is
optimistic by 4.9 percentage points. We consider the leakage-free number the
one to plan against.

\paragraph*{Model capacity, not the class ratio, was the binding constraint.}
Logistic Regression reached a precision-constrained recall of only 0.1479 on
exactly the same resampled data where Random Forest reached 0.7817. This
distinction reframes the standard advice: imbalance handling is necessary but
not sufficient, and the choice of classifier family dominates.

\paragraph*{Random Undersampling is actively harmful.} Its mean cost is 2.4
times that of SMOTE and 4.3 times the best cost achievable with SMOTE, and it
collapses precision-constrained recall to 0.063. Despite being among the most
commonly prescribed techniques, it is the worst choice on this dataset.

\paragraph*{The default importance measure is misleading.} Random Forest's
built-in impurity importance ranks \texttt{Amount} 21st of 35 features,
because the heavy-tailed amount distribution inflates impurity decrease.
Permutation importance and SHAP both rank it first or second. Had we reported
\texttt{feature\_importances\_}, we would have omitted the most informative
feature in the dataset.

\paragraph*{The operating policy matters more than the model.} A single
threshold cannot satisfy the 95\% precision floor and the 100-review-per-day
capacity at the same time: it either needs 1.9 analysts, or it must be
throttled to 40.9\% recall. The two-tier policy resolves the dilemma, reaching
81.0\% recall with only 70 reviews per day. The cheapest single threshold
near that operating point (0.1767) reaches a comparable recall of 80.3\% but
needs 232 reviews per day, which is more than twice the available capacity. The gain is therefore
operational rather than financial: the two-tier design buys the analyst-time
saving that the cost model alone does not price. This is the clearest path from
a research artefact to an operational deployment.

\subsection{Limitations}

\paragraph*{Only two days of data, from one region.} The dataset spans
approximately 48 hours of European card transactions. It cannot capture
seasonality, weekday/weekend cycles, or the drift that makes fraud patterns
change. A temporal split would have been preferable to a random split, but two
days cannot support one, so all temporal findings are indicative only.

\paragraph*{Anonymised PCA features limit interpretability.} \texttt{V1}--\texttt{V28}
are principal components; they cannot be given a business meaning. SHAP
quantifies their contribution but cannot explain \emph{why}. Our most important
feature, \texttt{Amount}, is interpretable; most of the model's power is not.

\paragraph*{One cost assumption carries substantial weight.} The
cost-minimising threshold moves with the review cost: it rises from 0.0405 at
USD 0.10 to 0.4605 at USD 2.00 (Table~\ref{tab:sens}). If the true cost of
blocking a legitimate customer is higher --- churn, card replacement,
reputational damage --- the bank should demand higher precision and accept
lower recall, and the 95\% precision constraint becomes the binding one
instead. We regard this as the most important sensitivity in the entire report,
and the model is explicitly parameterised so a bank can re-derive the optimum
for its own cost structure. Our recommended policy was selected under the
capacity constraint rather than the cost criterion, which makes it robust to
this assumption: auto-blocking consumes no analyst time at any review price.

\paragraph*{Sample size limits statistical power.} The test set contains only 142
frauds. A recall of 0.7817 has a 95\% confidence interval of roughly
$[0.71, 0.84]$, so differences below a few percentage points between
configurations should not be over-interpreted. This is also why the wider
hyper-parameter search did not pay off: with 331 training frauds, the search
space itself is a plausible source of overfitting, although --- as
Section~\ref{sec:search} sets out --- a single seed cannot establish that as a
proof.

\paragraph*{Our own thresholds are also test-derived.} We are explicit about
this because it cuts against our own conclusions. The operating thresholds of
the benchmark, of the single-threshold comparison \emph{and} of the
recommended two-tier policy were all selected on test-set predictions. The
model family, the hyper-parameters and the split itself were not. The
consequence is that the recall figures in this report are optimistic by an
amount we have quantified at 4.9 percentage points for the single-threshold
case, and that the two-tier thresholds (0.810 and 0.175) in particular must be
re-derived out-of-fold before they are used operationally. The qualitative
conclusions --- that Random Undersampling is harmful, that a tree ensemble is
required, that two tiers relax the capacity constraint --- do not depend on
the exact threshold values; the specific operating point does.

\paragraph*{Model selection on the test set.} Twenty-seven configurations were
compared on the same 142 test frauds, and the winner was then reported on
that same test set. This is a form of selection bias, and it is the reason we
present a tie explicitly in Section~\ref{sec:results}: two configurations
reach an identical operating point, and our choice among them rests on a
tie-break declared in advance (Average Precision) rather than on a
statistically significant difference. A dedicated validation set or nested cross-validation would remove
this bias; we did not reserve one, and with 331 training frauds the option was
not attractive.

\paragraph*{Unsupervised detection was not evaluated at scale.} Isolation Forest
was tested as a single unsupervised baseline. A hybrid pipeline that combines
supervised scoring with unsupervised anomaly flags for analyst attention is a
natural extension that we did not have the scope to implement.

\paragraph*{No deployment validation.} Latency was measured in-process on a
single machine; production characteristics such as throughput under concurrent
load, model serialisation size and monitoring of prediction drift remain
untested.

\subsection{Recommendations}

\begin{enumerate}
    \item \textbf{Deploy the two-tier policy.} Auto-block above 0.810, route
    the 0.175--0.810 band to analysts, approve the remainder. It attains higher
    recall than any single threshold that respects the capacity constraint,
    while requiring only 70 reviews per day, i.e.\ inside the capacity of one
    analyst.
    \item \textbf{Score every transaction.} Amount-based gating discards 90\% of
    the fraud. The model costs under USD 2.30 per year in compute at national
    scale, so there is no economic reason to restrict its input.
    \item \textbf{Prefer SMOTE or class weighting; avoid Random Undersampling.}
    Random Undersampling is the only technique that degraded every model
    family it was paired with, and its mean cost is 2.4 times that of SMOTE.
    Note that SMOTE gives the best ranking but not the lowest mean cost --- on
    average SMOTE-ENN is cheaper, because it trades recall for precision.
    \item \textbf{Plan capacity on the leakage-free numbers.} The test-tuned
    threshold of 0.52 and the two-tier thresholds reported here were all
    derived from test-set predictions and must be re-derived from out-of-fold
    training predictions before they are used for staffing or budget decisions;
    the leakage-free single-threshold recall is 73.24\% rather than 78.17\%.
    \item \textbf{Report permutation importance and SHAP, never
    \texttt{feature\_importances\_}} on heavy-tailed features.
    \item \textbf{Extend the dataset.} Merchant category, device fingerprint
    and transaction-history features are the missing ingredients for the
    low-value frauds that dominate our false negatives.
    \item \textbf{Monitor four quantities in production.} (i) The realised
    fraud rate in the auto-block tier --- a fall below the 0.980 precision we
    measured would mean the model is blocking legitimate customers and should
    trigger an immediate threshold review; (ii) the analyst queue depth
    against the 100-per-day capacity, with an alert at 80 transactions; (iii)
    the population stability index of the feature distributions, since
    \texttt{Amount} carries most of the model's signal and a shift in the
    transaction mix would invalidate the threshold; and (iv) the recall of the
    confirmed-fraud sample, estimated monthly, which is the only way to
    detect that the fraud pattern itself has drifted. Because the model is
    trained on a two-day window, we would also re-fit on a rolling 30-day
    window before any drift becomes material.
\end{enumerate}
